\documentclass[12pt,a4paper]{article}

%% ── Packages ────────────────────────────────────────────────────────────────
\usepackage[T1]{fontenc}
\usepackage{listings}
\usepackage{lstautogobble}
\usepackage[utf8]{inputenc}
\usepackage{lmodern}
\usepackage[margin=2.5cm]{geometry}
\usepackage{graphicx}
\usepackage{booktabs}
\usepackage{multirow}
\usepackage{amsmath}
\usepackage{amssymb}
\usepackage{caption}
\usepackage{subcaption}
\usepackage{hyperref}
\usepackage{xcolor}
\usepackage{array}
\usepackage{tabularx}
\usepackage{float}
\usepackage{enumitem}
\usepackage{microtype}
\usepackage{fancyhdr}
\usepackage{titlesec}
\usepackage{abstract}
\setlength{\headheight}{14pt}

%% ── Colours & Hyperref ──────────────────────────────────────────────────────
\definecolor{myblue}{RGB}{30, 80, 160}
\definecolor{mygreen}{RGB}{34, 120, 60}
\definecolor{mygray}{RGB}{90, 90, 90}
\hypersetup{
    colorlinks=true,
    linkcolor=myblue,
    urlcolor=myblue,
    citecolor=myblue,
    pdftitle={Boundary-Aware Residual Attention U-Net for Satellite LULC Segmentation},
    pdfauthor={},
}

%% ── Section title formatting ────────────────────────────────────────────────
\titleformat{\section}{\large\bfseries\color{myblue}}{{\thesection.}}{0.5em}{}[\titlerule]
\titleformat{\subsection}{\normalsize\bfseries}{{\thesubsection}}{0.5em}{}

%% ── Header / Footer ─────────────────────────────────────────────────────────
\pagestyle{fancy}
\fancyhf{}
\rhead{\small Boundary-Aware Residual Attention U-Net}
\lhead{\small Satellite LULC Segmentation — Sylhet, Bangladesh}
\cfoot{\thepage}
\renewcommand{\headrulewidth}{0.4pt}

%% ── Caption setup ───────────────────────────────────────────────────────────
\captionsetup{font=small, labelfont=bf, skip=4pt}

%% ─────────────────────────────────────────────────────────────────────────────
\begin{document}

%% ── Title Page ───────────────────────────────────────────────────────────────
\begin{titlepage}
  \centering
  \vspace*{2cm}
  {\LARGE\bfseries\color{myblue}
    Boundary-Aware Residual Attention U-Net\\[0.5em]
    for Satellite Image LULC Segmentation\par}
  \vspace{1.5cm}
  \rule{0.7\linewidth}{0.8pt}\\[1cm]
  {\large Project Report\\[0.4em]
   Deep Learning for Remote Sensing\par}
  \vspace{1.5cm}
  {\normalsize
    Study Area: Sylhet District, Bangladesh (2023)\\
    Sentinel-2 SR $\cdot$ ESRI LULC $\cdot$ 30\,m Resolution\\
    Trained on Kaggle GPU T4 $\cdot$ 200 Epochs per Run\\[1em]
    \today\par}
  \vfill
\end{titlepage}

\tableofcontents
\newpage

%% ─────────────────────────────────────────────────────────────────────────────
\section{Introduction and Motivation}

Land use and land cover (LULC) mapping from satellite imagery is a
fundamental task in environmental monitoring, urban planning, agricultural
management, and disaster response. Traditional methods rely on manual
digitization or spectral index thresholding, which are time-consuming and do
not scale to large areas.

This project builds an end-to-end deep learning pipeline for
\textbf{pixel-wise LULC classification} of the \textbf{Sylhet District,
Bangladesh} using a \textbf{Residual Attention U-Net} trained on
Sentinel-2 multispectral imagery paired with ESRI Global LULC ground-truth
labels.

The central research question is:

\begin{quote}
\itshape
Can we progressively improve satellite LULC segmentation — especially at
class transition boundaries — by (1)~fusing on-the-fly spectral indices
(NDVI, NDWI) into the input channels, and (2)~applying morphological
boundary-weighted loss?
\end{quote}

To answer this, a \textbf{5-run ablation study} was designed that isolates
the contribution of each component and, uniquely, sweeps the boundary loss
weight across three intensities ($\lambda \in \{0.5,\,1.0,\,2.0\}$) to
find the sweet spot between boundary sharpness and interior classification
quality.

%% ─────────────────────────────────────────────────────────────────────────────
\section{Study Area and Dataset}

\subsection{Study Area}

The study area is \textbf{Sylhet District, Bangladesh} — a geographically
diverse region encompassing haor wetland systems, tropical and sub-tropical
forest, dense tea-garden plantations, urban settlements, cultivated crop
fields, and exposed bare ground. This ecological variety makes Sylhet an
ideal benchmark for multi-class LULC segmentation: seven semantically
distinct classes must be discriminated within a single scene.

\subsection{Data Sources and Acquisition (Google Earth Engine)}

Both the satellite image and the LULC ground-truth mask were acquired,
composited, and exported using a custom \textbf{Google Earth Engine (GEE)}
JavaScript script.

\begin{table}[H]
\centering
\caption{Dataset summary.}
\begin{tabular}{ll}
\toprule
\textbf{Item} & \textbf{Details} \\
\midrule
Satellite Image Source & Sentinel-2 SR Harmonised (\texttt{COPERNICUS/S2\_SR\_HARMONISED}) \\
Label Source           & ESRI Global LULC 10\,m TS (\texttt{projects/sat-io/\ldots/ESRI\_Global-LULC\_10m\_TS}) \\
Study Area             & Sylhet District, Bangladesh \\
Year                   & 2023 (01 Jan -- 31 Dec) \\
Spatial Resolution     & 30\,m (exported via GEE) \\
Coordinate System      & UTM Zone 46N (EPSG:32646) \\
Sentinel-2 Bands       & B2 (Blue), B3 (Green), B4 (Red), B8 (NIR) \\
Cloud Filtering        & SCL-based pixel mask; $<$40\,\% cloud cover per scene \\
Composite Method       & Annual median of all valid pixels \\
Reflectance Scaling    & Divided by 10{,}000 to bring values into $[0,1]$ \\
\bottomrule
\end{tabular}
\end{table}

Clouds and cloud shadows were removed at pixel level using the Sentinel-2
\textbf{Scene Classification Layer (SCL)} before compositing. Valid pixels
were aggregated using the annual median, which suppresses remaining
atmospheric artefacts.

\subsection{LULC Label Preparation — 7-Class Remapping}

The ESRI LULC product defines ten class codes. \textbf{Snow/Ice} and
\textbf{Clouds} are entirely absent in the Sylhet study area and were
therefore excluded. The remaining seven classes were remapped to consecutive
integer labels 0--6. Any pixel originally carrying a Snow/Ice or Cloud code
was assigned the sentinel value 255, which is masked out of the loss function
and all metric calculations.

\begin{table}[H]
\centering
\caption{ESRI LULC codes remapped to 7 model labels.}
\begin{tabular}{cclc}
\toprule
\textbf{ESRI Code} & \textbf{Model ID} & \textbf{Class} & \textbf{Colour} \\
\midrule
1  & 0 & Water              & \texttt{\#1A5BAB} \\
2  & 1 & Trees              & \texttt{\#358221} \\
4  & 2 & Flooded Vegetation & \texttt{\#87D19E} \\
5  & 3 & Crops              & \texttt{\#FFDB5C} \\
7  & 4 & Built Area         & \texttt{\#ED022A} \\
8  & 5 & Bare Ground        & \texttt{\#EDE9E4} \\
9  & \textit{ignored (255)} & \textit{Snow/Ice}  & \textit{absent} \\
10 & \textit{ignored (255)} & \textit{Clouds}    & \textit{absent} \\
11 & 6 & Rangeland          & \texttt{\#C6AD8D} \\
\bottomrule
\end{tabular}
\end{table}

\begin{figure}[H]
  \centering
  \includegraphics[width=0.65\linewidth]{class_legend.png}
  \caption{Colour legend for the seven active LULC classes.}
\end{figure}

\subsection{Patch Extraction and Splitting}

The GeoTIFF images (approximately 370$\times$275 pixels at 30\,m) were
subdivided into $256 \times 256$ pixel patches using a sliding window with a
\textbf{step size of 128 pixels} (50\,\% overlap). This overlap strategy
quadruples the number of available training patches compared to
non-overlapping extraction, which is critical for a single-scene dataset.

Patches with fewer than 50\,\% finite (valid) Sentinel-2 pixels were
discarded. The surviving patches were randomly split 75\,\% train /
25\,\% validation using a \textbf{fixed random seed (42)} for
reproducibility. This yielded \textbf{243 training patches} and
\textbf{81 validation patches}.

%% ─────────────────────────────────────────────────────────────────────────────
\section{Model Architecture — Residual Attention U-Net}

The segmentation backbone is a \textbf{Residual Attention U-Net}: a
standard four-level U-Net augmented with two structural enhancements —
residual connections and attention gates.

\subsection{Encoder–Decoder Structure}

The encoder progressively halves the spatial resolution while doubling the
feature depth (64 $\to$ 128 $\to$ 256 $\to$ 512 filters), followed by a
1024-filter bottleneck. The decoder mirrors this with transposed convolutions
that restore spatial resolution level by level.

\subsection{Residual Connections}

Every convolutional block implements an \textbf{identity shortcut}:
\begin{equation}
  \mathbf{x}_\text{out} = \text{ReLU}\!\left(\mathcal{F}(\mathbf{x};\,\theta) + W_s\,\mathbf{x}\right)
\end{equation}
where $W_s$ is a $1\times1$ projection when channel dimensions differ.
Residual connections alleviate vanishing gradients and allow the network to
learn incremental corrections rather than full mappings.

\subsection{Attention Gates on Skip Connections}

Before each skip connection is concatenated with the decoder feature map, an
\textbf{attention gate} selectively suppresses irrelevant spatial activations:
\begin{align}
  \alpha &= \sigma\!\left(\mathbf{W}_a\,\text{ReLU}\!\left(\mathbf{W}_s\,\mathbf{s} + \mathbf{W}_g\,\mathbf{g}\right)\right) \\
  \hat{\mathbf{s}} &= \alpha \odot \mathbf{s}
\end{align}
where $\mathbf{s}$ is the skip feature from the encoder, $\mathbf{g}$ is the
gating signal from the decoder, and $\sigma$ is the sigmoid activation.
This allows the decoder to focus on task-relevant regions at each scale.

\subsection{Architecture Summary}

\begin{table}[H]
\centering
\caption{Layer configuration of the Residual Attention U-Net.}
\begin{tabular}{llr}
\toprule
\textbf{Block} & \textbf{Type} & \textbf{Filters} \\
\midrule
Encoder 1  & Residual Conv $\times$2 + MaxPool & 64 \\
Encoder 2  & Residual Conv $\times$2 + MaxPool & 128 \\
Encoder 3  & Residual Conv $\times$2 + MaxPool & 256 \\
Encoder 4  & Residual Conv $\times$2 + MaxPool & 512 \\
Bottleneck & Residual Conv $\times$2            & 1024 \\
Decoder 4  & TransConv + Attention Gate + Residual Conv $\times$2 & 512 \\
Decoder 3  & TransConv + Attention Gate + Residual Conv $\times$2 & 256 \\
Decoder 2  & TransConv + Attention Gate + Residual Conv $\times$2 & 128 \\
Decoder 1  & TransConv + Attention Gate + Residual Conv $\times$2 & 64 \\
Output     & Conv $1\times1$ + Softmax & 7 \\
\bottomrule
\end{tabular}
\end{table}

\begin{figure}[H]
  \centering
  \includegraphics[width=\linewidth]{unet_achtr.png}
  \caption{Architecture of the Residual Attention U-Net. The encoder
           (left) progressively halves the spatial resolution while
           doubling the feature depth (64$\to$128$\to$256$\to$512$\to$1024
           filters). Each skip connection passes through an
           \textbf{Attention Gate (AG)} that selectively weights encoder
           features using the decoder gating signal before concatenation.
           Every convolutional block includes a residual shortcut
           (bottom-right inset). The output head is a $1\times1$
           convolution with Softmax producing 7-class probability maps.}
\end{figure}

\noindent Dropout (rate\,$= 0.3$) is applied after every encoder block,
the bottleneck, and every decoder block to regularise training.

%% ─────────────────────────────────────────────────────────────────────────────
\section{Proposed Enhancements}

\subsection{Six-Channel Input via On-the-Fly Spectral Index Fusion}

The raw Sentinel-2 export provides four bands (B2=Blue, B3=Green, B4=Red,
B8=NIR). Two spectral indices are computed on-the-fly during data loading
and appended as additional input channels:

\begin{align}
  \text{NDVI} &= \frac{B8 - B4}{B8 + B4} \\[4pt]
  \text{NDWI} &= \frac{B3 - B8}{B3 + B8}
\end{align}

Both indices are clipped and linearly re-scaled to $[0,1]$ before
concatenation, expanding the input from 4 to \textbf{6 channels}. NDVI
discriminates photosynthetically active vegetation from soil and water;
NDWI highlights open water bodies and flooded areas. Including these
physics-derived features provides the network with domain knowledge that is
otherwise difficult to learn purely from raw reflectance values.

\subsection{Morphological Boundary-Weighted Focal Loss}

The total training loss is a 50/50 combination of \textbf{Dice loss} and
\textbf{Focal loss}. Class boundaries in satellite LULC maps are
inherently uncertain and often poorly delineated in 30\,m resolution data.
To guide the network to resolve these transitions more precisely, the
focal component is multiplied by a \textbf{morphological boundary weight
map}:

\begin{equation}
  B(L) = \text{dilate}(L,\,k) - \text{erode}(L,\,k), \quad k = 3\times3
\end{equation}

Pixels identified as lying on class boundaries receive an amplified loss:
\begin{equation}
  w_\text{boundary}(\mathbf{p}) = 1 + \lambda \cdot B(\mathbf{p})
\end{equation}

The full combined loss is:
\begin{equation}
  \mathcal{L} = 0.5\,\mathcal{L}_\text{Dice} + 0.5\,\mathcal{L}_\text{Focal}^\text{boundary}
\end{equation}

\subsection{Inverse-Frequency Class Weighting}

To handle the severe pixel-count imbalance across the seven classes,
per-class weights are computed as:
\begin{equation}
  w_c = \frac{1/\sqrt{f_c}}{\sum_j 1/\sqrt{f_j}}
\end{equation}
where $f_c$ is the fraction of training pixels belonging to class $c$.
This prevents dominant classes (Crops: $\approx$2.6\,M pixels) from
overwhelming rare ones (Bare Ground: $\approx$2.4\,K pixels).

%% ─────────────────────────────────────────────────────────────────────────────
\section{Training Configuration}

\begin{table}[H]
\centering
\caption{Hyperparameters and hardware.}
\begin{tabular}{ll}
\toprule
\textbf{Parameter} & \textbf{Value} \\
\midrule
Epochs              & 200 (early stopping, patience = 30) \\
Batch size          & 16 \\
Optimiser           & Adam \\
Learning rate       & $1\times10^{-4}$ (initial) \\
LR schedule         & Cosine Annealing \\
Loss function       & $0.5 \times \mathcal{L}_\text{Dice} + 0.5 \times \mathcal{L}_\text{Focal}^\text{boundary}$ \\
Dropout rate        & 0.3 \\
Global seed         & 42 (Python / NumPy / TensorFlow) \\
Hardware            & Kaggle GPU T4 ($\approx$32 min per run) \\
\bottomrule
\end{tabular}
\end{table}

\noindent A \textbf{global random seed of 42} was set via
\texttt{tf.keras.utils.set\_random\_seed(42)} prior to training. This seeds
Python, NumPy, and TensorFlow simultaneously, guaranteeing that weight
initialisation and batch shuffling are identical across all five ablation
runs. Without a fixed seed, stochastic variation in weight initialisation
on small datasets (243 training patches) can cause run-to-run swings of
3--5\% in Mean IoU, obscuring the true effect of each component.

%% ─────────────────────────────────────────────────────────────────────────────
\section{Evaluation Metrics}

Two complementary families of metrics were computed:

\subsubsection*{Pixel-Level Metrics}
\begin{itemize}
  \item \textbf{Overall Accuracy (OA):} fraction of correctly classified pixels.
  \item \textbf{Mean IoU (mIoU):} average Intersection-over-Union over all 7 classes.
  \item \textbf{Weighted IoU (wIoU):} IoU weighted by class pixel count.
  \item \textbf{Mean F1 Score:} harmonic mean of per-class precision and recall, averaged.
\end{itemize}

\subsubsection*{Boundary-Aware Metrics}
Boundary regions are extracted via the morphological gradient applied to the
\emph{predicted} and \emph{ground-truth} label maps:
\begin{equation}
  B_\text{pred} = \text{dilate}(P, k) - \text{erode}(P, k)
\end{equation}
\begin{itemize}
  \item \textbf{BF-Score (Boundary F1):} $2\cdot P_b \cdot R_b / (P_b + R_b)$, where $P_b$ and $R_b$ are boundary precision and recall.
  \item \textbf{Boundary IoU:} IoU computed exclusively on boundary pixels.
\end{itemize}

%% ─────────────────────────────────────────────────────────────────────────────
\section{Ablation Study Design}

Five configurations were evaluated sequentially to isolate the contribution
of each proposed component and to sweep the boundary loss multiplier:

\begin{table}[H]
\centering
\caption{Five ablation configurations.}
\begin{tabular}{clcccc}
\toprule
\textbf{Ver.} & \textbf{Description} & \textbf{Res.+Attn} & \textbf{NDVI/NDWI} & \textbf{Boundary Loss} & \textbf{$\lambda$} \\
\midrule
V1 & Baseline                   & \checkmark & $\times$   & $\times$   & --- \\
V2 & Baseline + NDVI/NDWI       & \checkmark & \checkmark & $\times$   & --- \\
V3 & Proposed (Soft boundary)   & \checkmark & \checkmark & \checkmark & 0.5 \\
V4 & Proposed (Medium boundary) & \checkmark & \checkmark & \checkmark & 1.0 \\
V5 & Proposed (Strong boundary) & \checkmark & \checkmark & \checkmark & 2.0 \\
\bottomrule
\end{tabular}
\end{table}

\noindent V1 serves as the baseline: it uses all architectural components
(residual connections + attention gates) but receives only 4-channel input and
no boundary weighting. Each subsequent version adds exactly one variable,
enabling clean attribution of performance changes.

%% ─────────────────────────────────────────────────────────────────────────────
\section{Results}

\subsection{Training Curves}

\begin{figure}[H]
  \centering
  \begin{subfigure}[b]{0.48\linewidth}
    \includegraphics[width=\linewidth]{v1_training_history.png}
    \caption{V1 Baseline}
  \end{subfigure}\hfill
  \begin{subfigure}[b]{0.48\linewidth}
    \includegraphics[width=\linewidth]{v3_training_history.png}
    \caption{V5 Proposed (Mult 2.0)}
  \end{subfigure}
  \caption{Training and validation loss/accuracy curves for Baseline (V1) vs.\ the full Proposed model (V5). Both runs converge smoothly without overfitting, indicating that the 50\,\% patch overlap provided sufficient training diversity.}
\end{figure}

\subsection{Global Performance Comparison}

\begin{table}[H]
\centering
\caption{Overall metrics for all five ablation versions.}
\begin{tabular}{lcccccc}
\toprule
\textbf{Metric} & \textbf{V1} & \textbf{V2} & \textbf{V3 (0.5)} & \textbf{V4 (1.0)} & \textbf{V5 (2.0)} & \textbf{Gain (V1$\to$V5)} \\
\midrule
Overall Accuracy    & 90.74\% & 91.19\% & 91.55\% & 91.52\% & \textbf{91.69\%} & $+0.95$\% \\
Mean IoU            & 70.14\% & 72.43\% & 73.53\% & 74.15\% & \textbf{74.98\%} & $+4.84$\% \\
Weighted IoU        & 83.81\% & 84.56\% & 85.14\% & 85.17\% & \textbf{85.40\%} & $+1.59$\% \\
Mean F1 Score       & 81.10\% & 82.91\% & 83.73\% & 84.18\% & \textbf{84.86\%} & $+3.76$\% \\
BF Score (boundary) & 0.6184  & 0.6326  & 0.6404  & 0.6493  & \textbf{0.6571}  & $+0.0387$ \\
Boundary IoU        & 44.76\% & 46.26\% & 47.10\% & 48.07\% & \textbf{48.93\%} & $+4.17$\% \\
Training Time (s)   & 1898    & 1907    & 1907    & 1916    & 1911             & ---       \\
\bottomrule
\end{tabular}
\end{table}

Every global metric improves \textbf{monotonically} from V1 through to V5,
with the single minor exception of Overall Accuracy between V3 and V4 which
differs by less than 0.04\,\% (well within rounding). This clean progression
confirms that both proposed contributions — spectral index fusion and
boundary-weighted loss — are genuinely additive.

\subsection{Per-Class IoU Comparison}

\begin{table}[H]
\centering
\caption{Per-class IoU (\%) across ablation versions.}
\begin{tabular}{lcccccc}
\toprule
\textbf{Class} & \textbf{V1} & \textbf{V2} & \textbf{V3 (0.5)} & \textbf{V4 (1.0)} & \textbf{V5 (2.0)} & \textbf{Net Gain} \\
\midrule
Water               & 94.01 & 94.22 & 94.30 & 94.46 & \textbf{94.49} & $+0.48$\% \\
Trees               & 69.85 & 70.84 & 72.52 & 71.64 & \textbf{72.61} & $+2.76$\% \\
Flooded Vegetation  & 46.00 & 50.96 & 52.54 & 54.08 & \textbf{55.53} & $+9.53$\% \\
Crops               & 88.35 & 88.62 & 88.98 & 89.08 & \textbf{89.13} & $+0.78$\% \\
Built Area          & 55.73 & 59.41 & 61.05 & 61.09 & \textbf{61.60} & $+5.87$\% \\
Bare Ground         & 85.60 & 88.53 & 89.53 & \textbf{91.37} & 90.32 & $+4.72$\% \\
Rangeland           & 51.42 & 54.44 & 55.76 & 57.29 & \textbf{61.20} & $+9.78$\% \\
\bottomrule
\end{tabular}
\end{table}

\begin{figure}[H]
  \centering
  \begin{subfigure}[b]{0.48\linewidth}
    \includegraphics[width=\linewidth]{v1_per_class_iou.png}
    \caption{V1 Baseline}
  \end{subfigure}\hfill
  \begin{subfigure}[b]{0.48\linewidth}
    \includegraphics[width=\linewidth]{v3_per_class_iou.png}
    \caption{V5 Proposed (Mult 2.0)}
  \end{subfigure}
  \caption{Per-class IoU bar charts for Baseline vs.\ Proposed model.}
\end{figure}

\subsection{Confusion Matrices}

\begin{figure}[H]
  \centering
  \begin{subfigure}[b]{0.48\linewidth}
    \includegraphics[width=\linewidth]{v1_confusion_matrix.png}
    \caption{V1 Baseline}
  \end{subfigure}\hfill
  \begin{subfigure}[b]{0.48\linewidth}
    \includegraphics[width=\linewidth]{v3_confusion_matrix.png}
    \caption{V5 Proposed (Mult 2.0)}
  \end{subfigure}
  \caption{Normalized confusion matrices. Off-diagonal improvements are most visible for Rangeland (row 6) and Flooded Vegetation (row 2).}
\end{figure}

\subsection{All-Metrics Chart — Proposed Model}

\begin{figure}[H]
  \centering
  \includegraphics[width=0.9\linewidth]{v3_all_metrics.png}
  \caption{Grouped bar chart of IoU, Precision, Recall, and F1 for each class in V5 (Proposed).}
\end{figure}

\subsection{Proposed Model — Per-Class Boundary Metrics}

\begin{table}[H]
\centering
\caption{Boundary metrics for V5 (Proposed -- Mult 2.0) per class.}
\begin{tabular}{lcccc}
\toprule
\textbf{Class} & \textbf{Boundary Prec.} & \textbf{Boundary Rec.} & \textbf{BF Score} & \textbf{Boundary IoU} \\
\midrule
Water              & 72.7\% & 53.0\% & 0.613 & 0.442 \\
Trees              & 66.8\% & 63.5\% & 0.651 & 0.483 \\
Flooded Vegetation & 72.9\% & 31.9\% & 0.444 & 0.285 \\
Crops              & 65.1\% & 57.4\% & 0.609 & 0.439 \\
Built Area         & 56.9\% & 42.1\% & 0.484 & 0.319 \\
Bare Ground        & \textbf{98.3\%} & \textbf{87.3\%} & \textbf{0.925} & \textbf{0.860} \\
Rangeland          & 81.4\% & 41.7\% & 0.552 & 0.381 \\
\bottomrule
\end{tabular}
\end{table}

\begin{figure}[H]
  \centering
  \includegraphics[width=0.9\linewidth]{v3_boundary_metrics.png}
  \caption{Per-class boundary metrics chart for V5 (Proposed).}
\end{figure}

\subsection{Prediction Visualizations}

\begin{figure}[H]
  \centering
  \includegraphics[width=\linewidth]{v3_predictions.png}
  \caption{Sample predictions from V5 (Proposed). Columns left to right: RGB Input, NDVI Index, Ground Truth, Model Prediction.}
\end{figure}

\subsection{Boundary Prediction Visualizations}

\begin{figure}[H]
  \centering
  \includegraphics[width=\linewidth]{v3_boundary_predictions.png}
  \caption{Boundary prediction overlay for V5 (Proposed).
           \textcolor{green}{\textbf{Green}} = True Positive boundary (correctly predicted),
           \textcolor{red}{\textbf{Red}} = False Negative (missed boundary),
           \textcolor{yellow}{\textbf{Yellow}} = False Positive (spurious boundary).}
\end{figure}

%% ─────────────────────────────────────────────────────────────────────────────
\section{Discussion and Key Findings}

\begin{enumerate}[leftmargin=*, label=\textbf{\arabic*.}]

  \item \textbf{Spectral index fusion (V1$\to$V2) yields the single largest
        per-step jump.}
        Adding NDVI and NDWI channels contributed $+2.29$\% Mean IoU and
        improved every class except Bare Ground, which remained stable.
        Flooded Vegetation gained $+4.96$\% — consistent with NDWI's
        sensitivity to inundated areas.

  \item \textbf{Boundary-weighted loss is monotonically beneficial as
        $\lambda$ increases.}
        Moving from no boundary loss (V2) to $\lambda=2.0$ (V5) added a
        further $+2.55$\% Mean IoU and $+2.45$\% BF Score. Importantly,
        increasing $\lambda$ never degraded overall accuracy; this is in
        contrast to earlier (unseeded) experiments where random
        initialisation caused the strong boundary penalty to destabilise
        rare-class learning.

  \item \textbf{Rangeland is the biggest beneficiary of boundary
        weighting} ($+9.78$\% IoU, V1$\to$V5). Rangeland spectrally
        overlaps with Crops and Trees; forcing the network to resolve
        their shared transitions helps it learn a more discriminative
        internal representation for the class.

  \item \textbf{Flooded Vegetation remains the most challenging class}
        (55.53\% IoU at best) due to its spectral similarity with both
        open water and dense vegetation, and because it is the rarest
        class ($\approx$52\,K pixels vs.\ $\approx$2.6\,M for Crops).

  \item \textbf{Bare Ground peaks at V4 ($\lambda=1.0$, 91.37\%) and
        drops slightly at V5 ($\lambda=2.0$, 90.32\%).} Although V5 wins
        overall, Bare Ground prefers the medium boundary weight, likely
        because its hard-edged, compact patches are already easy to
        delineate and the stronger boundary loss introduces marginal
        over-constraint.

  \item \textbf{Boundary recall is the bottleneck.} In the per-class
        boundary table the model consistently achieves higher boundary
        precision than recall, meaning it produces fewer spurious
        boundaries than it misses true ones. Future improvements could
        focus on recall — e.g.\ through a larger morphological kernel,
        a dedicated boundary detection head, or multi-scale boundary
        supervision.

  \item \textbf{Dominant classes generalise strongly.}
        Water (94.49\%), Crops (89.13\%), and Bare Ground (90.32\%) all
        achieve top performance, validating the inverse-frequency class
        weighting strategy.

  \item \textbf{Seeding was essential for reliable ablation.}
        Early unseeded runs showed run-to-run variance of up to
        $\pm$16\% IoU on Rangeland, masking the true signal. Fixing the
        global seed (Python / NumPy / TensorFlow all set to 42) produced
        clean, reproducible monotonic trends across all five runs.

\end{enumerate}

%% ─────────────────────────────────────────────────────────────────────────────
\section{Limitations and Future Work}

\begin{table}[H]
\centering
\caption{Current limitations and suggested improvements.}
\begin{tabularx}{\linewidth}{lX}
\toprule
\textbf{Limitation} & \textbf{Suggested Improvement} \\
\midrule
Single-scene dataset (one GeoTIFF tile) & Multi-temporal or multi-region training for better generalisation \\
30\,m resolution limits fine-boundary accuracy & Use 10\,m Sentinel-2 bands (B2, B3, B4, B8 at native resolution) \\
Boundary recall lag & Add a dedicated CASENet-style boundary detection head \\
No data augmentation & Rotation, flipping, colour jitter, and Mixup for robustness \\
Label noise at class transitions & Soft-label or confidence-weighted supervision near boundaries \\
\bottomrule
\end{tabularx}
\end{table}

%% ─────────────────────────────────────────────────────────────────────────────
\section{Conclusion}

This project demonstrates that a \textbf{Residual Attention U-Net} enhanced
with \textbf{on-the-fly spectral index fusion (NDVI/NDWI)} and a
\textbf{morphological boundary-weighted focal loss} achieves consistent,
monotonically improving performance over a strong attention-residual baseline
for satellite-image LULC segmentation of Sylhet, Bangladesh.

The 5-run ablation study cleanly isolates the contribution of each component
and shows that increasing the boundary loss weight from $\lambda = 0$ to
$\lambda = 2.0$ never sacrifices interior classification accuracy when the
training process is properly seeded and the dataset is generated with
sufficient patch overlap.

The proposed model (V5, $\lambda=2.0$) reaches:
\begin{itemize}
  \item \textbf{91.69\%} Overall Accuracy
  \item \textbf{74.98\%} Mean IoU ($+4.84$\% over the baseline)
  \item \textbf{BF Score 0.6571} ($+6.26$\% relative to baseline)
  \item \textbf{48.93\%} Boundary IoU ($+9.32$\% relative to baseline)
\end{itemize}

All five models were trained in under 32 minutes per run on a single Kaggle
T4 GPU, making the entire study reproducible at minimal computational cost.

%% ─────────────────────────────────────────────────────────────────────────────
\section*{Acknowledgements}
\begin{itemize}
  \item \textbf{Sentinel-2} — European Space Agency (ESA) multispectral satellite imagery.
  \item \textbf{ESRI LULC} — Esri 10-class global land use / land cover map (2023).
  \item \textbf{Google Earth Engine} — cloud-based geospatial data acquisition and preprocessing.
  \item \textbf{Kaggle} — GPU compute platform for model training.
\end{itemize}

\end{document}
